% !PRINTED 684
\weberpage{684}
% [source note] The same-parity restriction on x,y below applies only to odd m. For a fundamental discriminant -m the general condition is x congruent to m*y modulo 2; for m=4, x=0,y=1 represents i. The printed blanket restriction is retained.
Both determinants therefore have the minimal value, and
both systems are fundamental systems.

The case $\nu = 1$, which we excluded above, occurs, apart
from the case of the rational field, in which only the two
units $\pm 1$ exist, in the case of the \gesperrt{imaginary quadratic
 field}. In this case we have represented, in the twentieth
 chapter, the integers of the field in the form

\[
\frac{x + y\sqrt{-m}}{2}
\]

in which $-m$ denotes a fundamental discriminant, so
that $m \equiv 0$ or $\equiv 3\ (\mathrm{mod}\ 4)$, and $x, y$ are both even or
both odd. We obtain the units when we solve the
equation

\[
x^2 + m y^2 = 4
\]

in all possible ways in integral rational numbers. But this
equation has, when $m > 4$, only the two solutions
$x = \pm 2$, $y = 0$, and so in these cases, as in the
field of rational numbers, there are only the two units $\pm 1$. If
$m = 4$, one finds the four units $\pm 1, \pm i$, and if finally
 $m = 3$, the six units

\[
\pm 1, \; \pm \frac{-1 + i\sqrt{3}}{2}, \; \pm \frac{-1 - i\sqrt{3}}{2}.
\]

In the next simplest case, $n = 2$, $\nu = 2$, that is, in the
real quadratic field, the theory of units coincides
with the theory of the Pell equation treated in \S.\,127 of the
first volume.

\begin{center}
\S.\,188.\\[4pt]
\gesperrt{Reduced Numbers.}
\end{center}

If $\alpha$ is any integral or fractional number of the field $\Omega$
 different from zero, then we consider, as with the
units, the \gesperrt{system of conjugate logarithms} of $\alpha$,
by which we understand, as in \S.\,185, the \gesperrt{real numbers}

\weq{1}{\lambda_1 = \delta_1 \log|\alpha_1|, \; \lambda_2 = \delta_2 \log|\alpha_2|, \; \ldots, \; \lambda_\nu = \delta_\nu \log|\alpha_\nu|}

which we also denote by $\lambda_s(\alpha)$. If we consult the
fundamental system $\varepsilon_1, \varepsilon_2, \ldots, \varepsilon_{\nu-1}$ of units, then
%CONT>>
